\section*{ВВЕДЕНИЕ}
\addcontentsline{toc}{section}{ВВЕДЕНИЕ}

Интеллектуальные системы распознавания на основе изображений стали неотъемлемой частью современных технологий, стремительно развиваясь с момента появления первых алгоритмов компьютерного зрения. Первые шаги в этой области были связаны с простыми задачами классификации и детекции объектов, однако сегодня такие системы способны решать сложные задачи: от медицинской диагностики до автономного управления транспортом. Интенсивность развития технологий обработки изображений не имеет аналогов — они трансформируют промышленность, медицину, безопасность и повседневную жизнь. Современные системы видеонаблюдения, беспилотные автомобили, автоматизированные производственные линии и даже социальные сети невозможно представить без алгоритмов машинного зрения.

Интеллектуальные системы распознавания основаны на методах глубокого обучения и нейронных сетей, которые позволяют анализировать изображения с высокой точностью, в отличие от традиционных алгоритмов, требующих ручной настройки признаков. Благодаря этому они способны адаптироваться к изменяющимся условиям, улучшая свою производительность с увеличением объема данных.

Современные компании, осознавая потенциал искусственного интеллекта, активно внедряют системы компьютерного зрения для оптимизации бизнес-процессов. Например, в розничной торговле такие технологии помогают анализировать поведение покупателей, в промышленности — контролировать качество продукции, а в медицине — автоматизировать диагностику заболеваний. Разработка интеллектуальных систем распознавания позволяет не только повысить эффективность работы предприятий, но и создать новые сервисы, привлекающие клиентов.

Сайт или мобильное приложение, использующее технологии компьютерного зрения, может стать мощным инструментом для привлечения аудитории. Например, системы дополненной реальности (AR) или сервисы распознавания объектов в реальном времени способны значительно улучшить пользовательский опыт. Поэтому создание интеллектуальной системы на основе обработки изображений — это не только техническая задача, но и возможность вывести бизнес на новый уровень.

\emph{Цель настоящей работы} – разработка интеллектуальной системы, объединяющей методы компьютерного зрения и машинного обучения для распознавания лиц и отпечатков пальцев, интегрированной с картами доступа. Для достижения поставленной цели необходимо решить \emph{следующие задачи:}
\begin{itemize}
\item провести анализ предметной области и существующих решений;
\item разработать архитектуру системы и выбрать методы распознавания;
\item спроектировать настольное приложение для распознавания на основе изображений лиц и отпечатков пальцев;
\item реализовать приложение, используя графический интерфейс.
\end{itemize}

\emph{Структура и объем работы.} Отчет состоит из введения, 4 разделов основной части, заключения, списка использованных источников, 2 приложений. Текст выпускной квалификационной работы равен \formbytotal{lastpage}{страниц}{е}{ам}{ам}.

\emph{Во введении} сформулирована цель работы, поставлены задачи разработки, описана структура работы, приведено краткое содержание каждого из разделов.

\emph{В первом разделе} на стадии анализа предметной области изучается история развития систем распознавания, сравниваются существующие технологии распознавания и их преимущества и недостатки.

\emph{Во втором разделе} на стадии технического задания приводятся требования к разрабатываемой интеллектуальной системе распознавания.

\emph{В третьем разделе} на стадии технического проектирования представлены проектные решения для системы распознавания.

\emph{В четвертом разделе} приводится список классов и их методов, использованных при разработке системы распознавания, производится тестирование разработанного настольного приложения.

В заключении излагаются основные результаты работы, полученные в ходе разработки.

В приложении А представлен графический материал.
В приложении Б представлены фрагменты исходного кода. 
